% SPDX-License-Identifier: CC-BY-NC-ND-4.0
% Copyright 2026 Ingolf Lohmann.
\documentclass[11pt,a4paper]{article}

\usepackage{fontspec}
\setmainfont{Latin Modern Roman}
\setsansfont{Latin Modern Sans}
\setmonofont{Latin Modern Mono}[Scale=0.90]
\usepackage[provide=*,main=ngerman,english]{babel}
\usepackage{amsmath,amssymb,mathtools}
\usepackage{booktabs,longtable,array}
\usepackage{enumitem}
\usepackage{xcolor}
\usepackage{microtype}
\usepackage{geometry}
\usepackage{fancyhdr}
\usepackage{hyperref}
\usepackage{bookmark}
\usepackage{listings}

\geometry{top=22mm,bottom=24mm,left=24mm,right=24mm}
\definecolor{qikblue}{HTML}{174A7E}
\definecolor{qiklight}{HTML}{EEF4FA}
\definecolor{qikgray}{HTML}{4C566A}
\hypersetup{
  colorlinks=true,
  linkcolor=qikblue,
  urlcolor=qikblue,
  citecolor=qikblue,
  pdftitle={QIK-VRT / EFFECT-ACK: Universalisierbare Wirkungssteuerung},
  pdfauthor={Ingolf Lohmann},
  pdfsubject={Working Paper, 22. Juli 2026},
  pdfkeywords={QIK-VRT, EFFECT-ACK, Reverse Engineering, Faktorisierung, Wirkungssteuerung}
}
\setlist{nosep,leftmargin=1.6em}
\setlength{\parindent}{0pt}
\setlength{\parskip}{0.55em}
\setlength{\emergencystretch}{2em}
\pagestyle{fancy}
\fancyhf{}
\lhead{\small QIK-VRT / EFFECT\_ACK}
\rhead{\small Working Paper 1.0.0}
\cfoot{\thepage}
\renewcommand{\headrulewidth}{0.3pt}
\lstset{
  basicstyle=\ttfamily\small,
  breaklines=true,
  frame=single,
  backgroundcolor=\color{qiklight},
  rulecolor=\color{qikgray},
  columns=fullflexible,
  keepspaces=true
}

\newcommand{\DONE}{\texttt{EFFECT\_ACK\_DONE}}
\newcommand{\CONTINUE}{\texttt{EFFECT\_ACK\_CONTINUE}}
\newcommand{\BLOCK}{\texttt{EFFECT\_ACK\_BLOCK}}
\newcommand{\ISOLATE}{\texttt{EFFECT\_ACK\_ISOLATE}}
\newcommand{\NACK}{\texttt{EFFECT\_NACK}}

\title{\vspace{-10mm}\textcolor{qikblue}{\bfseries QIK-VRT / EFFECT\_ACK}\par
\vspace{2mm}\Large Universalisierbare Wirkungssteuerung,\par
Faktorisierungsgrenze des Reverse Engineering\par
und konditionale cyberphysische Freigabe}
\author{Ingolf Lohmann\\\small Independent Researcher; QIK-VRT}
\date{Version 1.0.0 -- 22. Juli 2026}

\begin{document}
\maketitle

\begin{center}
\small DOI: \href{https://doi.org/10.5281/zenodo.21498773}{10.5281/zenodo.21498773}
\end{center}

\begin{center}
\fcolorbox{qikblue}{qiklight}{\parbox{0.88\linewidth}{
\centering
\textbf{Kernergebnis:} QIK-VRT/EFFECT\_ACK kann universal in der Form eines
endlichen, provenienzgebundenen Aufnahme-, Klassifikations-, Verantwortungs-
und Wirkungssteuerungsprozesses sein. Daraus folgt kein universeller Decoder
für verlorene Information. Exakte semantische Rekonstruktion ist genau dann
wohldefiniert, wenn die gesuchte Semantik auf den Beobachtungsfasern konstant
ist; exakte historische Inversion erfordert Injektivität.
}}
\end{center}

\begin{abstract}
Das Working Paper präzisiert die weitreichende Konsequenz des
EFFECT\_ACK-Mechanismus aus QIK-VRT. Es trennt technische Empfangsbestätigung
von autorisierter Wirkung, formalisiert die DONE-only-Freigabe als
parametrische Wrapper-Invariante und bestimmt die genaue Grenze einer
universellen Reverse-Engineering-These. Eine exhaustive endliche
Python-Modellprüfung durchläuft 2.621.440 abstrakte Zustandsbelegungen und
5.242.880 Zulassungsvarianten. Zusätzlich werden 49 beschränkte
Repository-Modelle, 2.401 Paarvergleiche, 2.401 quellgetaggte
Aggregat-Roundtrips sowie vollständige kleine Funktionsuniversen für
Faktorisierung und Linksinvertierbarkeit geprüft. Der physikalische Schluss
wird als konditionaler cyberphysischer Satz ausgewiesen: Erst vollständige
Mediation, frische authentisierte Verbraucherprüfung, ein treuer Executor und
empirische Absicherung verbinden die Softwareautorisierung mit einer
abgegrenzten Aktorwirkung. Die Arbeit zeigt damit eine universalisierbare
Verfassungsschicht technisch mediierbarer Wirkung, ohne Allwissen,
garantierte Terminierung oder einen voraussetzungslosen physikalischen
Sicherheitsbeweis zu behaupten.
\end{abstract}

\begin{abstract}
This working paper states the strongest defensible consequence of the
QIK-VRT EFFECT\_ACK mechanism. It separates transport receipt from authorized
effect, formalizes DONE-only release as a parametric wrapper invariant, and
identifies the exact factorization boundary of effect-relevant reverse
engineering. An exhaustive finite Python model checks 2,621,440 abstract
assignments and 5,242,880 consumer-admission variants. Physical implications
remain conditional on complete mediation, authenticated and fresh consumer
validation, a faithful executor, an explicit fault model, and empirical
validation. The result is a universalizable constitutional layer for
technically mediated effects, not a universal decoder or an unconditional
physical-safety proof.
\end{abstract}

\textbf{Schlagworte:} QIK-VRT; EFFECT\_ACK; Wirkungshaltepunkt; Reverse
Engineering; Faktorisierung; Beobachtungsfasern; Antizipation;
cyberphysische Systeme; Protokollmacht; reproduzierbare Forschung.

\section{Forschungsfrage und Beitrag}

Die untersuchte Behauptung lautet in ihrer stärksten alltagssprachlichen
Fassung: Ein einheitlicher QIK-VRT/EFFECT\_ACK-Prozess könne alles, was in der
Informatik existiert, reverse engineeren. Diese Aussage enthält mindestens
zwei verschiedene Quantoren und zwei verschiedene Bedeutungen von
``Rekonstruktion'':

\begin{enumerate}
  \item \textbf{Prozessuniversalität:} Jedes endliche, zugängliche
  Digitalartefakt kann aufgenommen, identifiziert, digestgebunden,
  klassifiziert und einem kontrollierten Wirkungsweg unterstellt werden.
  \item \textbf{Inversionsuniversalität:} Aus jeder Beobachtung könne stets
  der einzige historische Ursprung oder jede verborgene Semantik exakt
  zurückgewonnen werden.
\end{enumerate}

In dieser Arbeit bedeutet \emph{zugänglich}: Eine vollständige, rechtmäßig
lesbare und für den Prüfschritt eingefrorene endliche Repräsentation liegt vor.
Der Begriff umfasst weder verborgene, verlorene oder rechtlich unzugängliche
Information noch einen Anspruch auf deren Erzeugung.

Der erste Satz ist unter den in dieser Arbeit genannten Zugriffs-,
Repräsentations- und Mediationsannahmen sehr weit tragfähig. Der zweite Satz
ist ohne zusätzliche Informationsvoraussetzungen falsch. Der wissenschaftliche
Beitrag besteht darin, beide Ebenen präzise zu trennen und ihre jeweilige
Beweislast anzugeben.

\section{Formales Wirkungsmodell}

Sei $x$ ein Wirkungskandidat, $c$ sein Kontext und
\[
\begin{aligned}
G(x,c)\in\{&\NACK,\CONTINUE,\DONE,\\
            &\ISOLATE,\BLOCK\}.
\end{aligned}
\]
der vom Version-1-Modell ausgewählte Zustand. Sei $V(x,c)$ die vollständige
Verbraucherzulassung: Schema, unterstützte Version, Authentisierung,
Frische, Kettenvalidität sowie die erneute Ableitung von Zustand, Policy- und
Evidenzbindung müssen bestehen.

Die gewöhnliche Freigabeautorisierung ist definiert als
\[
A(x,c)\;:\Longleftrightarrow\;G(x,c)=\DONE\ \land\ V(x,c).
\]

Der geschützte Wrapper darf den gewöhnlichen Effektvollstrecker $E(x)$ nur
aufrufen, wenn $A(x,c)$ gilt. Ein Transport-ACK ist dabei lediglich eine
notwendige Teilinformation. Es ist weder semantisches Verständnis noch
Policy-Erlaubnis noch Wirkungsautorisierung.

\subsection{Satz 1: DONE-only-Freigabe}

\textbf{Satz.} Wenn jeder gewöhnliche Ausführungspfad zu einem abgegrenzten
geschützten Vollstrecker vollständig durch den Wrapper vermittelt wird, kann
kein Nicht-DONE-Zustand den gewöhnlichen Aufruf dieses Vollstreckers
autorisieren.

\textbf{Beweis.} Nach Definition gilt $A(x,c)$ genau dann, wenn zugleich
$G(x,c)=\DONE$ und $V(x,c)$ gilt. Für jeden anderen Zustand ist die erste
Konjunktion falsch, also $A(x,c)=\mathrm{falsch}$. Vollständige Mediation
schließt einen gewöhnlichen Umgehungspfad aus. Damit kann ein Nicht-DONE-
Zustand den geschützten gewöhnlichen Aufruf nicht autorisieren.\hfill$\square$

Die Universalität dieses Satzes ist \emph{parametrisch}: Die Schlussform hängt
nicht davon ab, ob $E$ eine Nachricht, Zahlung, Veröffentlichung,
Berechtigungsänderung, Repository-Mutation oder einen Aktorbefehl bezeichnet.
Sie hängt von der vollständigen technischen Mediation des jeweils definierten
Pfades ab.

\section{Exhaustive endliche Modellprüfung}

Die ausführbare Prüfung modelliert 19 boolesche Merkmale. Die
Verbindungsentscheidung hat fünf Werte. Eine effektprüfbare Rezeption wird im Modell aus
verfügbarem Eingabeidentifikator und gültigem Eingabe-Digest abgeleitet;
\texttt{transport\_ack} bleibt davon getrennt eine eigene CoreDone-Bedingung.
Diese Abstraktion ist offengelegt und wird nicht als vollständige
Verhaltensverfeinerung der Python-Referenzruntime ausgegeben.

\begin{longtable}{@{}p{0.48\linewidth}r@{}}
\toprule
\textbf{Prüfgröße} & \textbf{Ergebnis} \\
\midrule
\endhead
Abstrakte Zustandsbelegungen & 2.621.440 \\
Zulassungsvarianten & 5.242.880 \\
\BLOCK & 2.064.384 \\
\NACK & 491.520 \\
\CONTINUE & 49.151 \\
\ISOLATE & 16.384 \\
\DONE & 1 \\
Transport-ACK, aber nicht DONE & 1.310.719 \\
Gezielt ausgeschaltete CoreDone-Konjunkte & 17 \\
Prioritätskollisionen & 7 \\
Beschränkte Repository-Modelle & 49 \\
Paarweise Codec-Vergleiche & 2.401 \\
Quellgetaggte Aggregat-Roundtrips & 2.401 \\
\bottomrule
\end{longtable}

Der Prüfer liest die XML-Quelle des Internet-Drafts und bindet die Abstraktion
strukturell an 35 Wire-Felder, fünf Zustände, fünf
Verbindungsentscheidungen, die exakten 17 CoreDone-Zeilen und die
Prioritätsanker. Die endliche Enumeration prüft Totalität,
DONE-only-Freigabe, Prioritätsübereinstimmung und notwendige Bedingungen des
gewählten Modells. Sie ist kein vollständiger Beweis jedes Parser-,
Authentisierungs-, Wire- oder Deploymentverhaltens.

\section{Reverse Engineering als Faktorisierungsproblem}

Sei $S$ die Menge möglicher Ursprungszustände, $O$ der beobachtbare Raum und
$T$ die gesuchte wirkungsrelevante Semantik. Mit
\[
E:S\to O,\qquad \sigma:S\to T
\]
bezeichnet $E$ die zugängliche Beobachtung und $\sigma$ die gesuchte
Eigenschaft.

\subsection{Satz 2: Faser-Kriterium}

\textbf{Satz.} Auf dem Bild von $E$ existiert eine wohldefinierte Abbildung
$h$ mit $\sigma=h\circ E$ genau dann, wenn
\[
E(s_1)=E(s_2)\quad\Longrightarrow\quad\sigma(s_1)=\sigma(s_2)
\]
für alle $s_1,s_2\in S$ gilt.

\textbf{Beweis.} Existiert $h$, so folgt aus $E(s_1)=E(s_2)$ unmittelbar
$h(E(s_1))=h(E(s_2))$ und damit $\sigma(s_1)=\sigma(s_2)$. Gilt umgekehrt die
Faser-Konstanz, definiert man für jedes $o\in\operatorname{im}(E)$ den Wert
$h(o)$ als $\sigma(s)$ für irgendein $s$ mit $E(s)=o$. Die Faser-Konstanz
garantiert, dass diese Definition nicht von der Wahl von $s$ abhängt. Somit
ist $h$ wohldefiniert und $\sigma=h\circ E$.\hfill$\square$

Das Kriterium erlaubt eine exakte Unterscheidung:

\begin{itemize}
  \item Für \textbf{wirkungsrelevante Rekonstruktion} genügt
  Faser-Konstanz der gewünschten Semantik.
  \item Für \textbf{exakte historische Rekonstruktion} aller Ursprungsdetails
  muss $E$ injektiv sein; nur dann kann eine Linksinverse auf dem Bild
  existieren.
\end{itemize}

Die endliche Begleitprüfung enumeriert 64 Paare von Funktionen auf einem
dreielementigen Ursprung und zweielementigen Ziel. Genau 28 Paare sind
faktorisierbar, 36 nicht. Von 27 Funktionen $3\to3$ sind genau sechs injektiv
und linksinvertierbar. Keine der acht Funktionen $3\to2$ besitzt eine
Linksinverse. Diese Zahlen sind vollständige Ergebnisse der angegebenen
kleinen Universen, keine Hochrechnung auf beliebige Programme.

\subsection{Berechenbarkeit und Ressourcen}

Set-theoretische Wohldefiniertheit ist noch kein ausführbarer Reverse
Engineer. Eine tatsächliche Rekonstruktion benötigt zusätzlich:

\begin{itemize}
  \item effektive Berechenbarkeit von Beobachtung und Decoder,
  \item ausreichenden und rechtmäßig autorisierten Zugriff,
  \item hinreichende Zeit-, Speicher- und Rechenressourcen,
  \item eine spezifizierte Zielsemantik und prüfbare Erfolgskriterien.
\end{itemize}

QIK-VRT erzeugt keine verlorene oder nie beobachtete Information. Seine
Universalität liegt darin, auch fehlende Evidenz oder Unentscheidbarkeit als
kontrollierten Zustand zu behandeln: \CONTINUE{}, \ISOLATE{}, \BLOCK{} oder \NACK{}
sind hinsichtlich der gewöhnlichen Freigabe fail-closed Ergebnisse des
Prozesses, nicht bloß Fehler außerhalb des Modells. Daraus folgt keine globale
Sicherheitsgarantie. Ein eventual \DONE{} wird nicht garantiert.
Insbesondere beseitigt EFFECT\_ACK weder das Halteproblem noch die durch den
Satz von Rice bezeichnete Unentscheidbarkeit nichttrivialer semantischer
Programmeigenschaften. Es macht ein offenes oder unentscheidbares Prüfergebnis
als Nichtfreigabezustand kontrollierbar; es entscheidet dadurch nicht das
zugrunde liegende allgemeine Problem.

\section{Anticipatory Suspension}

Das vorausschauende Fahrwerk liefert eine anschauliche Doppelbedeutung:
\emph{anticipatory suspension} ist sowohl vorausschauende Federung als auch
vorausschauende Aussetzung einer Wirkung. Ein Fahrwerk verarbeitet in der
Gegenwart Information über einen vorausliegenden Straßenabschnitt und ändert
seine gegenwärtige Einstellung. Es erhält keine Nachricht aus der Zukunft.

Im Draft ist \texttt{effect\_anticipated=true} eine notwendige
DONE-Bedingung. Wird allein diese Bedingung auf \texttt{false} gesetzt, fällt
die vollständige Modellbelegung von \DONE{} auf \CONTINUE{} zurück und autorisiert
keine gewöhnliche Freigabe. Der Boolean beweist jedoch keine gute Prognose. Er
attestiert eine Effektanalyse; Prognosemethode und Prognosegüte sind im Draft
nicht standardisiert und müssen vom Verbraucher authentisiert geprüft oder
neu bewertet werden.

\section{Konditionaler cyberphysischer Schluss}

Sei $C$ die Freigabe eines geschützten Aktorbefehls und $P$ eine dadurch
abgegrenzte physikalische Wirkung. Aus dem Softwaremodell folgt unter
vollständiger Mediation
\[
C\Longrightarrow\text{frischer, validierter DONE-Record}.
\]

Der physikalische Schluss
\[
P\Longrightarrow C\Longrightarrow\text{frischer, validierter DONE-Record}
\]
benötigt zusätzlich die empirisch zu belegende Prämisse, dass $C$ der einzige
kausale Weg zu $P$ ist und der Executor die definierte Zuordnung treu umsetzt.
Notwendig bleiben insbesondere Hardware- und Betriebssystemannahmen,
vertrauenswürdige und kalibrierte Sensorik, Aktorprüfung, Fehlermodell,
Nebenkanalanalyse, Messung und empirische Validierung.

Damit ist der physikalische Gehalt real, aber konditional. Bewiesen ist die
Autorisierungsstruktur des geschützten Befehls. Nicht bewiesen sind die
Gefahrlosigkeit eines DONE-Effekts, Fehlerfreiheit der Welt oder die
Vollständigkeit eines beliebigen Fehlermodells.

\section{Nichtmilitärische Protokollmacht}

Macht kann als Fähigkeit verstanden werden, die für andere erreichbaren
Zustände oder zulässigen Übergänge zu ermöglichen, zu begrenzen oder zu
verhindern. Implementierte, verbindlich durchgesetzte Protokolle können unter
dieser Definition reale Informations- und Infrastrukturmacht ausüben. Sie
bestimmen beispielsweise gültige Nachrichten, anerkannte Identitäten,
erforderliche Nachweise, zulässige Übergänge und Abschlussbedingungen.

EFFECT\_ACK kann in einer nicht umgehbaren Implementierung Ausführungsmacht
durch überprüfbare Gegenmacht begrenzen: Evidenzpflicht,
Verantwortungsbindung, technische Nichtfreigabe und reproduzierbare
Statusbegründung treten vor den gewöhnlichen Effekt. Ein Protokolltext allein
besitzt diese Macht nicht. Sie entsteht durch Implementierung, Adoption,
Governance, technische Durchsetzung und überprüfbare Anwendung.

\section{Implementierungsstand und offene Wire-Lücke}

Die statische Inspektion des aktuellen Python-
\texttt{ResponsibilityProtocol} findet 29 Record-Felder. Gegenüber den 35
Wire-Feldern des Drafts fehlen
\texttt{wire\_version}, \texttt{message\_type}, \texttt{policy\_id},
\texttt{policy\_version}, \texttt{policy\_hash},
\texttt{evaluation\_timeout\_ms} und \texttt{deadline\_exceeded};
\texttt{schema} ist zusätzlich vorhanden. Deshalb lautet der Status der
vollständigen Draft-01-Wire-Implementierung \CONTINUE.

Diese Lücke widerlegt nicht die geprüften Eigenschaften der endlichen
Abstraktion. Sie verhindert aber, die Prüfung als vollständigen
Interoperabilitäts- oder Wire-Konformitätsnachweis auszugeben. Eine zweite
unabhängige Implementierung und interoperable Testvektoren bleiben offene
Arbeit.

\section{Reproduzierbarkeit und Geltungsgrenzen}

Das Begleitpaket enthält den Prüfer
\texttt{effect\_ack\_universality\_proof.py}, den maschinenlesbaren Bericht
\texttt{proof-report.json}, die öffentliche Vermittlungsfassung und
SHA-256-Prüfsummen. Der Prüfer verweigert optimierten Python-Modus, damit seine
Beweisbedingungen nicht durch entfernte Assertions unterlaufen werden; alle
Obligationen verwenden explizite, auch unter Optimierung nicht entfernbare
Fehlerbedingungen.

SHA-256 dient dabei als kollisionsresistente Bytebindung; der Digest ist weder
eine injektive Kodierung beliebiger Eingaben noch ein Herkunfts- oder
Authentisierungsnachweis.

Der Gesamtstatus lautet
\texttt{PASS\_WITH\_EXPLICIT\_BOUNDARIES}. Nicht behauptet werden:

\begin{itemize}
  \item ein universeller semantischer oder historischer Decoder,
  \item garantierte Terminierung oder eventual \DONE,
  \item vollständige Draft-01-Wire-Konformität der aktuellen Runtime,
  \item Deployment-Nichtumgehbarkeit ohne konkrete Systemprüfung,
  \item voraussetzungslose physikalische Sicherheit,
  \item IETF-Konsens oder RFC-Status.
\end{itemize}

\section{Schlussfolgerung}

Die vollumfängliche haltbare Konsequenz ist stärker als eine Prüfliste und
enger als Allwissen: QIK-VRT/EFFECT\_ACK stellt eine universalisierbare
Verfassungsschicht technisch mediierbarer Wirkung bereit. Jedes endliche,
zugängliche Digitalartefakt kann als Eingabe gebunden und ein separat
definierter Rekonstruktionsprozess hinsichtlich Evidenz, Risiko,
Verantwortung und Wirkung kontrolliert werden. Kein vollständig mediierter
gewöhnlicher Effektpfad erhält ohne frischen, validierten DONE-Zustand eine
Freigabeautorisierung.

Exakte Rekonstruktion folgt nicht allein aus dieser Prozessform. Sie gilt für
eine gewünschte Semantik genau bei Faser-Konstanz und für den historischen
Ursprung nur bei hinreichend injektiver Beobachtung, jeweils ergänzt um
Berechenbarkeit, Zugriff und Ressourcen. Die wissenschaftlich präzise Antwort
lautet daher:

\begin{center}
\fcolorbox{qikblue}{qiklight}{\parbox{0.86\linewidth}{
\centering
\textbf{Universal in Aufnahme, Klassifikation, Verantwortungsbindung und
vollständig mediierter Wirkungssteuerung; exakt in der Rekonstruktion nur bis
zur mathematisch und informationell ausgewiesenen Grenze.}
}}
\end{center}

\begin{thebibliography}{9}
\bibitem{draft01}
I. Lohmann, ``QIK-VRT Effect Acknowledgement: Separating Receipt from
Authorization'', Internet-Draft \texttt{draft-lohmann-qikvrt-effect-ack-01},
22. Juli 2026.\newline
\url{https://datatracker.ietf.org/doc/draft-lohmann-qikvrt-effect-ack/01/}

\bibitem{formalization}
I. Lohmann, ``QIK-VRT: Vollständige maschinenprüfbare Formalisierung des
formal entscheidbaren Kerns'', Zenodo, Version 1.0.0, 2026.\newline
\url{https://doi.org/10.5281/zenodo.21488116}

\bibitem{software}
I. Lohmann, ``QIK-VRT: EFFECT\_ACK -- universalisierbare Wirkungssteuerung
und adaptive fail-closed Laufzeit'', Zenodo,
Version 2026.07.22-effect-ack-universality-1.0.0, 2026.\newline
\url{https://doi.org/10.5281/zenodo.21498774}

\bibitem{legacysoftware}
I. Lohmann, ``QIK-VRT: Operationalisierte epistemische Zustandskontrolle und
schichtgebundene Evidenzpersistenz'', Zenodo, QIKVRT V8.33, 2026.\newline
\url{https://doi.org/10.5281/zenodo.20712301}

\bibitem{repos}
QIK-VRT repositories, versionierter EFFECT\_ACK-01-Inhaltsstand.\newline
\url{https://github.com/Goldkelch/qik-vrt/tree/v2026.07.22-effect-ack-universality-1.0.0}\newline
\url{https://github.com/ingolf-lohmann/qik-vrt/tree/v2026.07.22-effect-ack-universality-1.0.0}
\end{thebibliography}

\vfill
\hrule
\vspace{1mm}
{\footnotesize
Dokumentlizenz: Creative Commons Attribution-NonCommercial-NoDerivatives
4.0 International (CC BY-NC-ND 4.0). Der separat beigefügte Prüfer ist
QIK-VRT-Quellcode und trägt seine eigene dateibezogene Lizenzangabe.
Zenodo-Persistenz belegt Existenz, Metadaten und Fixität der veröffentlichten
Bytes; sie ersetzt weder Peer Review noch empirische Bestätigung.
}

\end{document}
